\originalchapter{边染色}\label{ch:edgecolor}
\sourceof{18.315 L6--7; Vizing 扇形的旋转与换色为编者补证}
\originalsection{匹配分解}
\begin{definition}
正常边染色给每条边分配一种颜色, 使有共同端点的两边异色. 最少颜色数为边染色数 $\chi'(G)$. 每个边颜色类就是一个匹配, 所以边染色等价于把全部边分拆成匹配.
\end{definition}
在排班模型中, 顶点可以表示队伍, 边表示两队之间的比赛. 同一队同一时段最多参加一场, 一种边颜色就表示一个时段. 最大度给出必要时段数 $\chi'\ge\Delta$, 但奇圈会需要额外一色: $\chi'(C_{2r+1})=3$, 而 $\Delta=2$.
\begin{theorem}[二部图边染色]\label{thm:bipartiteedge}
有限二部多重图满足 $\chi'(G)=\Delta(G)$.
\end{theorem}
\begin{proof}
若 $\Delta=0$, 无边而结论成立. 在两侧补孤立顶点使它们大小相等. 令 $D=\Delta$, 每个点的缺额度为 $D-d(v)$. 两侧缺额度总和相等, 因为原两侧度数和都等于边数. 反复连接两侧尚有缺额的点, 允许平行边, 直到变成 $D$ 正则二部多重图. 由推论\ref{cor:regularmatching}有完美匹配. 删去它后图为 $D-1$ 正则, 归纳分拆成 $D$ 个完美匹配. 限制到原边, 得 $D$ 色边染色. 下界来自最大度点.
\end{proof}
这证明了 Hall 与边染色之间的联系: 完美匹配移走每个顶点的一单位度, 因而可以逐层剥离正则二部图. 一般图的奇圈使这层结构失效.
\originalsection{双色链与扇形}
以下只讨论简单图. 在一张正常部分边染色中, 只看颜色 $\alpha,\beta$ 的边, 每个顶点至多关联这两色各一条边, 所以分支为路、偶圈或孤立点. 在某个分支交换两色, 不会破坏正常边染色, 这称为 Kempe 换色.
\begin{lemma}[线性扇形旋转]\label{lem:fan}
设只有边 $xy_0$ 未染色, 邻点 $y_0,y_1,\ldots,y_t$ 互异, 且边 $xy_i$ 的颜色 $c_i$ 在 $y_{i-1}$ 处缺失 ($1\le i\le t$). 对任意前缀 $y_0,\ldots,y_s$, 同时把 $xy_{i+1}$ 的颜色移到 $xy_i$ ($i<s$), 让 $xy_s$ 未染, 仍是正常部分边染色.
\end{lemma}
\begin{proof}
中心 $x$ 处保留原有互异颜色的一部分, 没有重复. 在每个 $y_i$ ($i<s$) 处, 移入颜色原本缺失; 最后一个顶点只删去一种颜色. 其余边不变, 所以无冲突. 若颜色 $\gamma$ 在 $x$ 与 $y_s$ 同时缺失, 就可在旋转后给 $xy_s$ 染 $\gamma$.
\end{proof}
旋转必须以原边颜色同时重新赋值, 或从前到后依次正确移位; 不能误把刚改过的颜色继续当作原色传播.
\originalsection{Vizing 定理}
\begin{theorem}[Vizing]\label{thm:vizing}
有限简单图满足 $\Delta\le\chi'\le\Delta+1$.
\end{theorem}
\begin{proof}
固定 $D+1$ 色, 其中 $D=\Delta$. 对边数归纳, 设除 $xy_0$ 外全部边已正常染色. 每点至少缺一色. 固定 $x$ 处缺色 $\alpha$, 从单点扇形 $y_0$ 开始. 在当前末点 $y_t$ 取缺色 $\beta$. 若 $x$ 也缺 $\beta$, 旋转整个扇形后染末边 $\beta$, 完成. 否则 $x$ 恰有一条 $\beta$ 色边 $xz$. 若 $z$ 不在扇形中, 将它作为新末点继续. 中心邻点有限, 所以终会直接完成或重复.

若重复为 $z=y_j$, 则 $1\le j\le t-1$: 首边未染, 且 $y_t$ 缺 $\beta$, 所以重复边既不是首边也不是末边. 有 $c_j=\beta$, 而 $y_{j-1},y_t$ 都缺 $\beta$. 考察 $\alpha,\beta$ 双色子图中含 $x$ 的分支 $P$. $x$ 缺 $\alpha$ 而有 $\beta$, 所以 $P$ 是以 $x$ 为端点的路. 两个缺 $\beta$ 的点若在此路上, 只能成为另一端点, 因而 $y_{j-1},y_t$ 至多一个属于 $P$. 交换 $P$ 上的两色, 此后 $x$ 缺 $\beta$.

如果 $y_{j-1}\notin P$, 它仍缺 $\beta$. 前缀 $y_0,\ldots,y_{j-1}$ 的中心边原来都不是 $\alpha,\beta$: $\alpha$ 在 $x$ 缺失, 唯一 $\beta$ 边是未纳入该前缀的 $xy_j$. 前缀旋转涉及的其他缺色条件不受双色交换影响. 旋转此前缀, 再给其末边染 $\beta$, 完成.

如果 $y_{j-1}\in P$, 则 $y_t\notin P$. 换色后 $y_t$ 仍缺 $\beta$, $y_{j-1}$ 改为缺 $\alpha$, 而 $xy_j$ 由 $\beta$ 改为 $\alpha$. 因而扇形在这个位置的新条件仍成立: 新边色 $\alpha$ 在前点缺失. 其余中心边都不是这两色, 其缺色条件不变. 所以可旋转整个扇形, 再把末边染 $\beta$. 归纳完成. 下界仍由最大度点的关联边必须异色得到.
\end{proof}
这里没有把“沿双色路换色”作为未经验证的口号. 两种情况分别控制换色是否改动旋转前缀, 这正是原手写稿略写、必须补齐的地方. 简单图条件也很重要: 平行边会让不同中心边共用另一端点, 上述不同邻点扇形构造不能直接使用.
\originalsection{完全图的时段安排}
\begin{proposition}
$\chi'(K_{2r})=2r-1$, $\chi'(K_{2r+1})=2r+1$ ($r\ge1$).
\end{proposition}
\begin{proof}
对 $K_{2r}$, 把顶点标为 $\infty$ 及模 $2r-1$ 的剩余类. 对每个 $a$, 取匹配
\[
\{\infty a\}\ \cup\ \{(a+i)(a-i):1\le i\le r-1\}.
\]
模数为奇数, 非零剩余类可按 $\pm i$ 成对, 所以每个匹配覆盖全部顶点. 任意两有限标号 $u,v$ 的边归于唯一中点 $a=(u+v)/2$, 除2在奇模数下可逆; $\infty$ 的边也各出现一次. 得到 $2r-1$ 色, 与最大度下界相等.

$K_{2r+1}$ 每个匹配最多含 $r$ 条边, 而总边数为 $r(2r+1)$, 所以至少需 $2r+1$ 色. 把它嵌入 $K_{2r+2}$, 用刚才的 $2r+1$ 色安排再删除额外顶点即可.
\end{proof}
\originalsection{习题与解答}
\begin{exercise}\label{ex:edge1}
证明图的边染色数等于其线图的顶点染色数. 线图的顶点为原图各边, 共享端点的边在新图中相邻.
\end{exercise}
\begin{exercise}\label{ex:edge2}
一所学校有若干老师与班级, 每条教学任务连接一位老师和一个班级. 每位老师及每个班级各参与至多 $D$ 条任务, 同一老师与班级可有多条任务. 证明能用 $D$ 个时段安排全部任务, 每时段老师和班级都没有冲突.
\end{exercise}
习题\ref{ex:edge1}的解答.
\begin{solution}
线图每个顶点颜色对应原边颜色. 线图相邻顶点异色恰好要求原图共端点两边异色, 两种染色逐一对应. 但在线图套一般 $\Delta+1$ 贪心界通常不如 Vizing, 因为原边 $uv$ 在线图度数最多 $d(u)+d(v)-2$.
\end{solution}
\Needspace{4\baselineskip}
习题\ref{ex:edge2}的解答.
\begin{solution}
任务构成二部多重图, 两侧为老师和班级. 二部图边染色定理给至多 $D$ 色, 每个颜色类为不共端点的匹配, 对应一个可同时执行的时段. 无需把重复任务合并, 因为多重图版本已保留它们.
\end{solution}
